\lecture{9.10}{Causal and Counterfactual Views of Missing Data}
\label{lec:9-10}

\subsection{Why Missingness Is Causal}

Missingness can be viewed as a causal process. Let $R_k$ indicate whether
variable $L_k$ is observed. Intervening to set $R_k=1$ corresponds to forcing the
variable to be measured. The full value can be written as a counterfactual
quantity:
\[
  L_k(1)=\text{the value of }L_k\text{ that would be recorded if }R_k=1.
\]
The factual observed value is
\[
  L_k=
  \begin{cases}
  L_k(1), & R_k=1,\\
  ?, & R_k=0.
  \end{cases}
\]
Thus missing data involve counterfactuals: values that would have been observed
under a measurement intervention.

\subsection{Missing Data DAGs}

A missing data DAG contains full-data variables, missingness indicators, and
observed coarsened variables. For one partially observed variable:

\begin{figure}[H]
\centering
\begin{tikzpicture}[
  box/.style={draw, rounded corners=2pt, minimum width=1.1cm, minimum height=.7cm},
  arr/.style={-{Latex[length=2.2mm]}, thick},
  det/.style={-{Latex[length=2.2mm]}, thick, gray},
  node distance=1.0cm
]
\node[box] (Lfull) {$L(1)$};
\node[box, below=of Lfull] (R) {$R$};
\node[box, right=of R] (Lobs) {$L$};
\draw[det] (Lfull) -- (Lobs);
\draw[det] (R) -- (Lobs);
\end{tikzpicture}
\caption{Basic missing-data structure: the observed value is determined by the full value and its missingness indicator.}
\label{fig:ch9-basic-missing-dag}
\end{figure}

With several variables, missingness indicators may depend on full values,
observed values, other indicators, or completely observed covariates $Z$.
Graphical restrictions encode assumptions about these dependencies.

\subsection{Full-Law Factorization}

Let $Z$ denote completely observed variables, $L(1)$ the full-data variables
that would be observed under measurement, $R$ the missingness indicators, and
$L$ the coarsened observed variables. The full law can be factored as
\[
  p\{Z,L(1),R,L\}
  =
  p\{Z,L(1)\}\,
  p\{R\given L(1),Z\}\,
  p\{L\given L(1),R,Z\}.
\]
The first factor is the target full-data law. The second factor is the
missingness mechanism. The third factor is deterministic once the consistency
rule is specified:
\[
  L_k=
  \begin{cases}
  L_k(1), & R_k=1,\\
  ?, & R_k=0.
  \end{cases}
\]
The observed-data law is obtained after marginalizing over the unobserved
components of $L(1)$.

\subsection{Recoverability}

\begin{definition}[Recoverability]
A target full-data parameter is recoverable if it can be written as a function
of the observed-data distribution.
\end{definition}

Recoverability is the missing-data analogue of causal identification. For
example, the full-data mean $\E\{L(1)\}$ is recoverable under suitable
missingness assumptions, but may not be recoverable if missingness depends on
the unobserved value in an unrestricted way.

\subsection{Missingness Mechanisms as Graphs}

Consider two partially observed variables, smoking status $L_1(1)$ and bronchitis
status $L_2(1)$, with indicators $R_1$ and $R_2$. Different arrows tell different
stories:
\begin{itemize}
  \item $L_2(1)\rightarrow R_1$: inquiry about smoking depends on true
  bronchitis status.
  \item $R_1\rightarrow R_2$: whether a diagnostic test is ordered depends on
  whether smoking was assessed.
  \item $L_1\rightarrow R_2$: observed smoking status affects diagnostic
  testing.
  \item $L_1(1)\rightarrow R_2$: true smoking status affects diagnostic testing,
  even when smoking status is not observed.
\end{itemize}

\begin{figure}[H]
\centering
\begin{tikzpicture}[
  box/.style={
    draw,
    rounded corners=2pt,
    align=center,
    minimum width=1.05cm,
    minimum height=.7cm
  },
  arr/.style={-{Latex[length=2.2mm]}, thick},
  det/.style={-{Latex[length=2.2mm]}, thick, gray}
]

% Nodes: symmetric layered layout
\node[box] (L11) at (0, 2.6) {$L_1(1)$};
\node[box] (L21) at (3.2, 2.6) {$L_2(1)$};

\node[box] (R1)  at (0, 1.25) {$R_1$};
\node[box] (R2)  at (3.2, 1.25) {$R_2$};

\node[box] (L1)  at (0, -0.15) {$L_1$};
\node[box] (L2)  at (3.2, -0.15) {$L_2$};

% Black structural arrows
\draw[arr] (L21.south west) to[out=210,in=35] (R1.north east);
\draw[arr] (R1.east) -- (R2.west);
\draw[arr] (L1.north east) to[out=35,in=220] (R2.south west);

% Gray deterministic arrows, routed outside
\draw[det] (L11.south west) to[out=-100,in=100] (L1.north west);
\draw[det] (R1.south west)  to[out=-120,in=60]  (L1.north east);

\draw[det] (L21.south east) to[out=-80,in=80]   (L2.north east);
\draw[det] (R2.south east)  to[out=-60,in=120]  (L2.north west);

\end{tikzpicture}
\caption{A missing-data DAG can encode how true values, observed values, and measurement decisions interact.}
\label{fig:ch9-missing-dag-example}
\end{figure}

\subsection{Graphical MCAR, MAR, and MNAR}

Graphically:
\begin{itemize}
  \item MCAR means there are no arrows into missingness indicators from data
  variables or other missingness-relevant variables.
  \item MAR means missingness indicators may depend on completely observed
  variables and observed values, but not on full values that are hidden when
  missing.
  \item MNAR means at least one missingness indicator has a parent among the
  full-data variables whose values may be unobserved.
\end{itemize}
The graph is more informative than a label because it records which variables
drive which missingness decisions.

\subsection{Workflow for Missing Data Analysis}

\begin{estimationbox}
A graphical missing-data workflow is:
\begin{enumerate}
  \item Define the full-data estimand, such as $\E\{L(1)\}$ or a causal effect.
  \item Draw a graph over full-data variables, missingness indicators, observed
  coarsened variables, and completely observed variables.
  \item Assess recoverability: can the estimand be written using the observed
  data law?
  \item If recoverable, estimate the identifying functional with regression,
  weighting, imputation, TMLE, or another appropriate method.
  \item If not recoverable, report bounds, perform sensitivity analysis, or
  collect additional data.
\end{enumerate}
\end{estimationbox}

\subsection{Connections to Censoring Bias}

Censoring is a special case of missingness over time. In survival analysis, a
future outcome is missing after dropout. Treating censoring indicators as
time-varying treatment variables leads to the same identification logic used for
longitudinal interventions: consistency, exchangeability for the censoring
process, and positivity of remaining uncensored.

\begin{keyidea}
The causal view treats missingness indicators as intervention variables and
missing values as counterfactual full-data quantities. This turns missing-data
analysis into a recoverability problem, parallel to causal identification.
\end{keyidea}
